\section{Conclusions}
\label{sec:conclusions}

CILP provides a feasible inverse-link-prediction formulation for utility-preserving sparsification of attributed social networks.
Across three attributed social networks---LastFM Asia, Facebook Page-Page, and GitHub Developers---each with ten paired seeds and budgets $0.1$--$0.9$, the multi-criteria teacher consistently improves sparsity--Macro-F1 AUC and structural metrics relative to an architecturally matched task-only teacher under identical fusion and exact budgets.
The fixed four-metric support rule is satisfied on all three datasets.
CILP remains competitive with ILP-GCN and PTDNet on predictive AUC, but it is not universally superior: on GitHub it significantly improves over Task-only yet remains slightly below ILP-GCN and PTDNet.
The supported methodological contribution is therefore composite counterfactual supervision with a controlled ablation, not overall sparsifier leadership.
CILP training, which includes teacher construction and scorer fitting, remains substantially more expensive than lighter baselines.
Future work includes training-only teacher targets, separate teacher/scorer timers, downstream GCN break-even measurement, leave-one-criterion-out and weight-sensitivity analyses, direct community-preservation metrics, full-graph ($0\%$) utility-retention references, adaptive teacher sampling, and validating the spectral proxy.
